% !PRINTED 385
\weberpage{385}
%<<CONT
und bezeichnen die hierdurch bedingte Permutation in doppelter Weise mit

\weq{3}{\Pi_{1,2,3,8} = \Pi_{4,5,6,7},}

indem wir festsetzen, dass $\Pi_{\alpha_1\alpha_2\alpha_3\alpha_4}$ und $\Pi_{\beta_1\beta_2\beta_3\beta_4}$ dasselbe bedeuten sollen, wenn $\beta_1,\beta_2,\beta_3,\beta_4$ irgend eine Permutation von $\alpha_1,\alpha_2,\alpha_3,\alpha_4$ ist, oder wenn $\alpha_1,\alpha_2,\alpha_3,\alpha_4,\beta_1,\beta_2,\beta_3,\beta_4$ zusammen alle acht Ziffern umfassen.

Durch dies Zeichen ist die Vertauschung (2) eindeutig bezeichnet, und die Anzahl der verschiedenen Permutationen dieser Art beträgt genau 35, so dass wir die ganze Gruppe $P$ durch die Nebengruppen so darstellen können:

\weq{4}{P = S + \Sigma\, S\, \Pi_{\alpha_1\alpha_2\alpha_3\alpha_4},}

worin $\alpha_1,\alpha_2,\alpha_3,\alpha_4$ alle Systeme von vier Ziffern aus der Reihe $1,2,\dots,8$ durchlaufen, wobei eine beliebige Ziffer, z.\,B. $\alpha_4 = 8$, festgehalten werden kann.

Um den Einfluss von $\Pi_{1,2,3,8}$ auf irgend eine Wurzel $[i\,k]$ zu erkennen, genügt es, die drei Wurzeln $[1\,2], [1\,4], [4\,5]$ zu betrachten, weil $1,2,3$ einerseits, $4,5,6,7$ andererseits ganz gleichartig in $\Pi_{1,2,3,8}$ vorkommen.

Diese Vertauschungen erhält man einfach aus der Bemerkung, dass $[1\,2], [1\,4]$ nach \S.\,99 die in dem Complex

\weq{5}{[2\,8][1\,2],\ [3\,8][1\,3],\ [4\,8][1\,4],\ [5\,8][1\,5],\ [6\,8][1\,6],\ [7\,8][1\,7]}

mit $[2\,8]$ und $[4\,8]$ verbundenen Wurzeln sind, und dass ebenso $[4\,5]$ die mit $[5\,8]$ verbundene Wurzel in dem Complexe

\weq{6}{[1\,8][1\,4],\ [2\,8][2\,4],\ [3\,8][3\,4],\ [5\,8][5\,4],\ [6\,8][6\,4],\ [7\,8][7\,4]}

ist. Durch die Vertauschung (2) gehen aber die Complexe (5) und (6) in folgende über, wie man leicht aus der Darstellung der Complexe im \S.\,99 findet [der Complex (5) bleibt als Ganzes ungeändert]:

\[
[3\,1][3\,8],\ [2\,1][2\,8],\ [4\,8][1\,4],\ [5\,8][5\,1],\ [6\,8][6\,1],\ [7\,8][7\,1]
\]

\[
[2\,3][1\,4],\ [3\,1][2\,4],\ [1\,2][3\,4],\ [5\,8][6\,7],\ [6\,8][5\,7],\ [7\,8][5\,6],
\]

woraus man folgende durch $\Pi_{1,2,3,8}$ bewirkte Vertauschungen erhält:

\weq{7}{\begin{array}{l} [1\,2],\ [1\,4],\ [4\,5] \\ {}[3\,8],\ [1\,4],\ [6\,7]. \end{array}}

Aus (2) und (7) lässt sich nun folgende allgemeine und einfache Regel ableiten.
